
In contemporary times, the LSM-Tree is widely adopted for its efficient insertion capabilities. However, a drawback arises when dealing with range deletes, as they are stored in individual SST files. Consequently, the process of searching for a deleted key necessitates opening and reading a file, leading to wasteful disk IO operations and time consumption. \\
This paper focuses on addressing this issue by proposing a range delete filter (RDF) to minimize disk IOs during the search for deleted keys. The RDF serves as a mechanism to swiftly identify and handle deleted keys, optimizing the overall performance of the LSM-Tree. \\